

An interaction occurs when the content of an item meets the taste of a
user, and the bilinear score of matrix-factorization models mirrors this
shape: a taste vector multiplied by a content vector~\cite{koren2009mf}.
The two ingredients differ in how well they are observed. The content of
an item is recorded; the catalogue states its department, colour, and
pattern. The taste of a user is latent. Users do not state their taste,
and in most cases could not articulate it if asked. The stated attributes
that are available, such as age or newsletter subscription, are a loose
proxy: sufficient signal for popularity-based fallbacks, but far too
unspecific to ground an explanation in. A recommendation justified with
``because you are 34 and receive the newsletter'' attributes the purchase
to facts that did not cause it.

Taste is nevertheless not out of reach. It cannot be collected as a
stated feature, but it is revealed through behavior~\cite{samuelson1948consumption}:
the purchase history records what a user's taste selects. Collaborative
filtering exploits exactly this, and the learned user embedding is a
summary of revealed taste. The user side therefore has no capture
problem. Its problem is that revealed taste arrives without labels, as a
latent vector with no vocabulary attached. For grounding, the consequence
is to stop searching for taste among user attributes and instead express
the captured taste in the vocabulary of the items it selects.

The same observation dictates where grounding must begin. In pure
collaborative filtering, both embedding spaces are behavior summaries;
the embedding of an item encodes who bought it, not what it is, and
whatever legibility the item side appeared to have was borrowed from the
catalogue vocabulary outside the model. The item side owns a rich
recorded vocabulary that can be moved inside the model, while the user
side offers only the thin stated fields discussed in
the Data chapter. Grounding therefore starts on the item side. The
order is forced by the data, and the user side follows once its revealed
taste can be expressed in item vocabulary (the fU chapter).
